% Options for packages loaded elsewhere
% Options for packages loaded elsewhere
\PassOptionsToPackage{unicode}{hyperref}
\PassOptionsToPackage{hyphens}{url}
\PassOptionsToPackage{dvipsnames,svgnames,x11names}{xcolor}
%
\documentclass[
  portuguese,
  letterpaper,
  DIV=11,
  numbers=noendperiod,
  oneside]{scrartcl}
\usepackage{xcolor}
\usepackage[left=1in,marginparwidth=2.0666666666667in,textwidth=4.1333333333333in,marginparsep=0.3in]{geometry}
\usepackage{amsmath,amssymb}
\setcounter{secnumdepth}{5}
\usepackage{iftex}
\ifPDFTeX
  \usepackage[T1]{fontenc}
  \usepackage[utf8]{inputenc}
  \usepackage{textcomp} % provide euro and other symbols
\else % if luatex or xetex
  \usepackage{unicode-math} % this also loads fontspec
  \defaultfontfeatures{Scale=MatchLowercase}
  \defaultfontfeatures[\rmfamily]{Ligatures=TeX,Scale=1}
\fi
\usepackage{lmodern}
\ifPDFTeX\else
  % xetex/luatex font selection
\fi
% Use upquote if available, for straight quotes in verbatim environments
\IfFileExists{upquote.sty}{\usepackage{upquote}}{}
\IfFileExists{microtype.sty}{% use microtype if available
  \usepackage[]{microtype}
  \UseMicrotypeSet[protrusion]{basicmath} % disable protrusion for tt fonts
}{}
\makeatletter
\@ifundefined{KOMAClassName}{% if non-KOMA class
  \IfFileExists{parskip.sty}{%
    \usepackage{parskip}
  }{% else
    \setlength{\parindent}{0pt}
    \setlength{\parskip}{6pt plus 2pt minus 1pt}}
}{% if KOMA class
  \KOMAoptions{parskip=half}}
\makeatother
% Make \paragraph and \subparagraph free-standing
\makeatletter
\ifx\paragraph\undefined\else
  \let\oldparagraph\paragraph
  \renewcommand{\paragraph}{
    \@ifstar
      \xxxParagraphStar
      \xxxParagraphNoStar
  }
  \newcommand{\xxxParagraphStar}[1]{\oldparagraph*{#1}\mbox{}}
  \newcommand{\xxxParagraphNoStar}[1]{\oldparagraph{#1}\mbox{}}
\fi
\ifx\subparagraph\undefined\else
  \let\oldsubparagraph\subparagraph
  \renewcommand{\subparagraph}{
    \@ifstar
      \xxxSubParagraphStar
      \xxxSubParagraphNoStar
  }
  \newcommand{\xxxSubParagraphStar}[1]{\oldsubparagraph*{#1}\mbox{}}
  \newcommand{\xxxSubParagraphNoStar}[1]{\oldsubparagraph{#1}\mbox{}}
\fi
\makeatother

\usepackage{color}
\usepackage{fancyvrb}
\newcommand{\VerbBar}{|}
\newcommand{\VERB}{\Verb[commandchars=\\\{\}]}
\DefineVerbatimEnvironment{Highlighting}{Verbatim}{commandchars=\\\{\}}
% Add ',fontsize=\small' for more characters per line
\usepackage{framed}
\definecolor{shadecolor}{RGB}{241,243,245}
\newenvironment{Shaded}{\begin{snugshade}}{\end{snugshade}}
\newcommand{\AlertTok}[1]{\textcolor[rgb]{0.68,0.00,0.00}{#1}}
\newcommand{\AnnotationTok}[1]{\textcolor[rgb]{0.37,0.37,0.37}{#1}}
\newcommand{\AttributeTok}[1]{\textcolor[rgb]{0.40,0.45,0.13}{#1}}
\newcommand{\BaseNTok}[1]{\textcolor[rgb]{0.68,0.00,0.00}{#1}}
\newcommand{\BuiltInTok}[1]{\textcolor[rgb]{0.00,0.23,0.31}{#1}}
\newcommand{\CharTok}[1]{\textcolor[rgb]{0.13,0.47,0.30}{#1}}
\newcommand{\CommentTok}[1]{\textcolor[rgb]{0.37,0.37,0.37}{#1}}
\newcommand{\CommentVarTok}[1]{\textcolor[rgb]{0.37,0.37,0.37}{\textit{#1}}}
\newcommand{\ConstantTok}[1]{\textcolor[rgb]{0.56,0.35,0.01}{#1}}
\newcommand{\ControlFlowTok}[1]{\textcolor[rgb]{0.00,0.23,0.31}{\textbf{#1}}}
\newcommand{\DataTypeTok}[1]{\textcolor[rgb]{0.68,0.00,0.00}{#1}}
\newcommand{\DecValTok}[1]{\textcolor[rgb]{0.68,0.00,0.00}{#1}}
\newcommand{\DocumentationTok}[1]{\textcolor[rgb]{0.37,0.37,0.37}{\textit{#1}}}
\newcommand{\ErrorTok}[1]{\textcolor[rgb]{0.68,0.00,0.00}{#1}}
\newcommand{\ExtensionTok}[1]{\textcolor[rgb]{0.00,0.23,0.31}{#1}}
\newcommand{\FloatTok}[1]{\textcolor[rgb]{0.68,0.00,0.00}{#1}}
\newcommand{\FunctionTok}[1]{\textcolor[rgb]{0.28,0.35,0.67}{#1}}
\newcommand{\ImportTok}[1]{\textcolor[rgb]{0.00,0.46,0.62}{#1}}
\newcommand{\InformationTok}[1]{\textcolor[rgb]{0.37,0.37,0.37}{#1}}
\newcommand{\KeywordTok}[1]{\textcolor[rgb]{0.00,0.23,0.31}{\textbf{#1}}}
\newcommand{\NormalTok}[1]{\textcolor[rgb]{0.00,0.23,0.31}{#1}}
\newcommand{\OperatorTok}[1]{\textcolor[rgb]{0.37,0.37,0.37}{#1}}
\newcommand{\OtherTok}[1]{\textcolor[rgb]{0.00,0.23,0.31}{#1}}
\newcommand{\PreprocessorTok}[1]{\textcolor[rgb]{0.68,0.00,0.00}{#1}}
\newcommand{\RegionMarkerTok}[1]{\textcolor[rgb]{0.00,0.23,0.31}{#1}}
\newcommand{\SpecialCharTok}[1]{\textcolor[rgb]{0.37,0.37,0.37}{#1}}
\newcommand{\SpecialStringTok}[1]{\textcolor[rgb]{0.13,0.47,0.30}{#1}}
\newcommand{\StringTok}[1]{\textcolor[rgb]{0.13,0.47,0.30}{#1}}
\newcommand{\VariableTok}[1]{\textcolor[rgb]{0.07,0.07,0.07}{#1}}
\newcommand{\VerbatimStringTok}[1]{\textcolor[rgb]{0.13,0.47,0.30}{#1}}
\newcommand{\WarningTok}[1]{\textcolor[rgb]{0.37,0.37,0.37}{\textit{#1}}}

\usepackage{longtable,booktabs,array}
\usepackage{calc} % for calculating minipage widths
% Correct order of tables after \paragraph or \subparagraph
\usepackage{etoolbox}
\makeatletter
\patchcmd\longtable{\par}{\if@noskipsec\mbox{}\fi\par}{}{}
\makeatother
% Allow footnotes in longtable head/foot
\IfFileExists{footnotehyper.sty}{\usepackage{footnotehyper}}{\usepackage{footnote}}
\makesavenoteenv{longtable}
\usepackage{graphicx}
\makeatletter
\newsavebox\pandoc@box
\newcommand*\pandocbounded[1]{% scales image to fit in text height/width
  \sbox\pandoc@box{#1}%
  \Gscale@div\@tempa{\textheight}{\dimexpr\ht\pandoc@box+\dp\pandoc@box\relax}%
  \Gscale@div\@tempb{\linewidth}{\wd\pandoc@box}%
  \ifdim\@tempb\p@<\@tempa\p@\let\@tempa\@tempb\fi% select the smaller of both
  \ifdim\@tempa\p@<\p@\scalebox{\@tempa}{\usebox\pandoc@box}%
  \else\usebox{\pandoc@box}%
  \fi%
}
% Set default figure placement to htbp
\def\fps@figure{htbp}
\makeatother



\ifLuaTeX
\usepackage[bidi=basic]{babel}
\else
\usepackage[bidi=default]{babel}
\fi
% get rid of language-specific shorthands (see #6817):
\let\LanguageShortHands\languageshorthands
\def\languageshorthands#1{}


\setlength{\emergencystretch}{3em} % prevent overfull lines

\providecommand{\tightlist}{%
  \setlength{\itemsep}{0pt}\setlength{\parskip}{0pt}}



 


\usepackage{booktabs}
\usepackage{longtable}
\usepackage{array}
\usepackage{multirow}
\usepackage{wrapfig}
\usepackage{float}
\usepackage{colortbl}
\usepackage{pdflscape}
\usepackage{tabu}
\usepackage{threeparttable}
\usepackage{threeparttablex}
\usepackage[normalem]{ulem}
\usepackage{makecell}
\usepackage{xcolor}
\KOMAoption{captions}{tableheading}
\makeatletter
\@ifpackageloaded{caption}{}{\usepackage{caption}}
\AtBeginDocument{%
\ifdefined\contentsname
  \renewcommand*\contentsname{Índice}
\else
  \newcommand\contentsname{Índice}
\fi
\ifdefined\listfigurename
  \renewcommand*\listfigurename{Lista de Figuras}
\else
  \newcommand\listfigurename{Lista de Figuras}
\fi
\ifdefined\listtablename
  \renewcommand*\listtablename{Lista de Tabelas}
\else
  \newcommand\listtablename{Lista de Tabelas}
\fi
\ifdefined\figurename
  \renewcommand*\figurename{Figura}
\else
  \newcommand\figurename{Figura}
\fi
\ifdefined\tablename
  \renewcommand*\tablename{Tabela}
\else
  \newcommand\tablename{Tabela}
\fi
}
\@ifpackageloaded{float}{}{\usepackage{float}}
\floatstyle{ruled}
\@ifundefined{c@chapter}{\newfloat{codelisting}{h}{lop}}{\newfloat{codelisting}{h}{lop}[chapter]}
\floatname{codelisting}{Listagem}
\newcommand*\listoflistings{\listof{codelisting}{Lista de Listagens}}
\makeatother
\makeatletter
\makeatother
\makeatletter
\@ifpackageloaded{caption}{}{\usepackage{caption}}
\@ifpackageloaded{subcaption}{}{\usepackage{subcaption}}
\makeatother
\makeatletter
\@ifpackageloaded{sidenotes}{}{\usepackage{sidenotes}}
\@ifpackageloaded{marginnote}{}{\usepackage{marginnote}}
\makeatother
\usepackage{bookmark}
\IfFileExists{xurl.sty}{\usepackage{xurl}}{} % add URL line breaks if available
\urlstyle{same}
\hypersetup{
  pdftitle={Anatomia de uma Bolha: O Ciclo do Subprime (1995-2015)},
  pdfauthor={Daniel Reis},
  pdflang={pt},
  colorlinks=true,
  linkcolor={blue},
  filecolor={Maroon},
  citecolor={Blue},
  urlcolor={Blue},
  pdfcreator={LaTeX via pandoc}}


\title{Anatomia de uma Bolha: O Ciclo do Subprime (1995-2015)}
\usepackage{etoolbox}
\makeatletter
\providecommand{\subtitle}[1]{% add subtitle to \maketitle
  \apptocmd{\@title}{\par {\large #1 \par}}{}{}
}
\makeatother
\subtitle{Uma autópsia econômica sobre expansão de crédito, CDOs e o
colapso do mercado}
\author{Daniel Reis}
\date{2026-03-05}
\begin{document}
\maketitle

\renewcommand*\contentsname{Índice}
{
\hypersetup{linkcolor=}
\setcounter{tocdepth}{3}
\tableofcontents
}

\begin{Shaded}
\begin{Highlighting}[]
\FunctionTok{source}\NormalTok{(}\StringTok{"\textasciitilde{}/Thesis/P1v2/preamble.R"}\NormalTok{)}
\end{Highlighting}
\end{Shaded}

\begin{verbatim}
Warning: package 'readr' was built under R version 4.2.3
\end{verbatim}

\begin{verbatim}
Warning: package 'readxl' was built under R version 4.2.3
\end{verbatim}

\begin{verbatim}
Warning: package 'MultipleBubbles' was built under R version 4.2.3
\end{verbatim}

\begin{verbatim}
Warning: package 'tidyverse' was built under R version 4.2.3
\end{verbatim}

\begin{verbatim}
Warning: package 'ggplot2' was built under R version 4.2.3
\end{verbatim}

\begin{verbatim}
Warning: package 'tibble' was built under R version 4.2.3
\end{verbatim}

\begin{verbatim}
Warning: package 'tidyr' was built under R version 4.2.3
\end{verbatim}

\begin{verbatim}
Warning: package 'purrr' was built under R version 4.2.3
\end{verbatim}

\begin{verbatim}
Warning: package 'dplyr' was built under R version 4.2.3
\end{verbatim}

\begin{verbatim}
Warning: package 'forcats' was built under R version 4.2.3
\end{verbatim}

\begin{verbatim}
Warning: package 'lubridate' was built under R version 4.2.3
\end{verbatim}

\begin{verbatim}
-- Attaching core tidyverse packages ------------------------ tidyverse 2.0.0 --
v dplyr     1.1.4     v purrr     1.0.2
v forcats   1.0.0     v stringr   1.5.0
v ggplot2   3.5.1     v tibble    3.2.1
v lubridate 1.9.3     v tidyr     1.3.1
-- Conflicts ------------------------------------------ tidyverse_conflicts() --
x dplyr::filter() masks stats::filter()
x dplyr::lag()    masks stats::lag()
x tibble::view()  masks seasonal::view()
i Use the conflicted package (<http://conflicted.r-lib.org/>) to force all conflicts to become errors
\end{verbatim}

\section{Introdução}\label{introduuxe7uxe3o}

Apesar de já terem se passado quase duas décadas, a crise (ou bolha)
conhecida como Crise do Subprime permanece como um importante fenômeno
econômico, financeiro e até mesmo comportamental, que merece ser
investigado e compreendido. O período que pretendemos descrever
refere-se ao crescimento explosivo dos preços imobiliários nos EUA (ao
menos em muitas regiões do país), à chegada ao seu valor máximo, até
então nunca visto, e ao subsequente colapso desses mesmos preços, sendo
assim, algo iniciado no mercado financeiro e imobiliário e que culminou
afetando a vida de boa parte do globo.

Muitos dos instrumentos financeiros, mecanismos de propagação, entre
outros fatores, permanecem sendo usados (embora não no mesmo contexto)
até os dias atuais. Este fatídico período da história econômica nos
rende, ademais, uma série de analogias de que podemos fazer uso
atualmente para compreender fenômenos econômicos e práticas financeiras
que persistem sendo observadas. Além disso, compreender o período nos
brinda com informações que nos auxiliem a diferenciar uma verdadeira
bolha especulativa de movimentos `naturais' dos mercados financeiros.

A ideia deste curto texto é, despretensiosamente, demonstrar da maneira
mais didática possível o que ocorreu naquele período: desde os fenômenos
precursores da bolha, quando talvez ela tenha de fato começado, quais
foram os combustíveis adicionados e a eclosão, até o retorno a uma ideia
de normalidade.

A tabela abaixo apresenta os principais eventos que se supõe terem
afetado o curso dos preços imobiliários nos EUA:

\begin{figure*}

\begin{Shaded}
\begin{Highlighting}[]
\FunctionTok{library}\NormalTok{(knitr)}
\FunctionTok{library}\NormalTok{(kableExtra)}
\end{Highlighting}
\end{Shaded}

\begin{verbatim}
Warning: package 'kableExtra' was built under R version 4.2.3
\end{verbatim}

\begin{verbatim}

Attaching package: 'kableExtra'
\end{verbatim}

\begin{verbatim}
The following object is masked from 'package:dplyr':

    group_rows
\end{verbatim}

\begin{Shaded}
\begin{Highlighting}[]
\FunctionTok{library}\NormalTok{(dplyr)}

\CommentTok{\# 1. Criar o dataframe de eventos com a linguagem técnica ajustada}
\NormalTok{tabela\_eventos }\OtherTok{\textless{}{-}} \FunctionTok{data.frame}\NormalTok{(}
  \AttributeTok{Data =} \FunctionTok{c}\NormalTok{(}\StringTok{"09/08/1995"}\NormalTok{, }\StringTok{"10/03/2000"}\NormalTok{, }\StringTok{"01/01/2001"}\NormalTok{, }\StringTok{"11/09/2001"}\NormalTok{, }\StringTok{"01/01/2004"}\NormalTok{, }
           \StringTok{"01/07/2006"}\NormalTok{, }\StringTok{"15/09/2008"}\NormalTok{, }\StringTok{"16/09/2008"}\NormalTok{, }\StringTok{"01/02/2012"}\NormalTok{),}
  \AttributeTok{Evento =} \FunctionTok{c}\NormalTok{(}\StringTok{"IPO da Netscape"}\NormalTok{, }\StringTok{"Pico do NASDAQ (Dot{-}com)"}\NormalTok{, }\StringTok{"Ciclo de Queda da Fed Funds Rate"}\NormalTok{, }
             \StringTok{"Ataques Terroristas de 11/09"}\NormalTok{, }\StringTok{"Mudança de Política das GSEs"}\NormalTok{, }
             \StringTok{"Pico de Preços (Housing Peak)"}\NormalTok{, }\StringTok{"Falência do Lehman Brothers"}\NormalTok{, }
             \StringTok{"Resgate da AIG"}\NormalTok{, }\StringTok{"Fundo do Mercado (Bottom)"}\NormalTok{),}
  \StringTok{\textasciigrave{}}\AttributeTok{Implicação}\StringTok{\textasciigrave{}} \OtherTok{=} \FunctionTok{c}\NormalTok{(}\StringTok{"Início da euforia tecnológica e semente da liquidez especulativa."}\NormalTok{,}
                 \StringTok{"Estouro da bolha de tecnologia; migração de capital para ativos reais."}\NormalTok{,}
                 \StringTok{"Início do afrouxamento monetário agressivo pelo Federal Reserve."}\NormalTok{,}
                 \StringTok{"Manutenção de juros em patamares baixos por período prolongado."}\NormalTok{,}
                 \StringTok{"Expansão da compra de hipotecas de baixa qualidade (Subprime)."}\NormalTok{,}
                 \StringTok{"Exaustão da demanda e início da correção nos preços nominais."}\NormalTok{,}
                 \StringTok{"Colapso da confiança e travamento do mercado interbancário."}\NormalTok{,}
                 \StringTok{"Reconhecimento do risco sistêmico e intervenção governamental."}\NormalTok{,}
                 \StringTok{"Estabilização dos preços após o ciclo de desalavancagem."}\NormalTok{)}
\NormalTok{)}

\CommentTok{\# 2. Gerar a tabela com kableExtra}
\NormalTok{tabela\_eventos }\SpecialCharTok{\%\textgreater{}\%}
  \FunctionTok{kbl}\NormalTok{(}\AttributeTok{caption =} \StringTok{"Cronologia de Eventos Críticos do Mercado Imobiliário (1995{-}2012)"}\NormalTok{,}
      \AttributeTok{booktabs =} \ConstantTok{TRUE}\NormalTok{, }
      \AttributeTok{align =} \FunctionTok{c}\NormalTok{(}\StringTok{"l"}\NormalTok{, }\StringTok{"l"}\NormalTok{, }\StringTok{"l"}\NormalTok{)) }\SpecialCharTok{\%\textgreater{}\%}
  \FunctionTok{kable\_styling}\NormalTok{(}\AttributeTok{bootstrap\_options =} \FunctionTok{c}\NormalTok{(}\StringTok{"striped"}\NormalTok{, }\StringTok{"hover"}\NormalTok{, }\StringTok{"condensed"}\NormalTok{),}
                \AttributeTok{full\_width =}\NormalTok{ F, }
                \AttributeTok{font\_size =} \DecValTok{14}\NormalTok{) }\SpecialCharTok{\%\textgreater{}\%}
  \FunctionTok{column\_spec}\NormalTok{(}\DecValTok{1}\NormalTok{, }\AttributeTok{bold =} \ConstantTok{TRUE}\NormalTok{)}
\end{Highlighting}
\end{Shaded}

\begin{table}
\centering
\caption{\label{tab:unnamed-chunk-2}Cronologia de Eventos Críticos do Mercado Imobiliário (1995-2012)}
\centering
\fontsize{14}{16}\selectfont
\begin{tabular}[t]{>{}lll}
\toprule
Data & Evento & Implicação\\
\midrule
\textbf{09/08/1995} & IPO da Netscape & Início da euforia tecnológica e semente da liquidez especulativa.\\
\textbf{10/03/2000} & Pico do NASDAQ (Dot-com) & Estouro da bolha de tecnologia; migração de capital para ativos reais.\\
\textbf{01/01/2001} & Ciclo de Queda da Fed Funds Rate & Início do afrouxamento monetário agressivo pelo Federal Reserve.\\
\textbf{11/09/2001} & Ataques Terroristas de 11/09 & Manutenção de juros em patamares baixos por período prolongado.\\
\textbf{01/01/2004} & Mudança de Política das GSEs & Expansão da compra de hipotecas de baixa qualidade (Subprime).\\
\addlinespace
\textbf{01/07/2006} & Pico de Preços (Housing Peak) & Exaustão da demanda e início da correção nos preços nominais.\\
\textbf{15/09/2008} & Falência do Lehman Brothers & Colapso da confiança e travamento do mercado interbancário.\\
\textbf{16/09/2008} & Resgate da AIG & Reconhecimento do risco sistêmico e intervenção governamental.\\
\textbf{01/02/2012} & Fundo do Mercado (Bottom) & Estabilização dos preços após o ciclo de desalavancagem.\\
\bottomrule
\end{tabular}
\end{table}

\end{figure*}%

\section{Capítulo 1: O ``Paciente Zero''
(1987)}\label{capuxedtulo-1-o-paciente-zero-1987}

A tecnologia da destruição é criada. O primeiro \textbf{CDO} nasce no
Drexel Burnham Lambert. * \textbf{O Conceito:} Transformar ``lixo''
(Junk Bonds) em títulos com selo AAA através de tranches. * \textbf{A
Semente:} A engenharia financeira prova que é possível ``esconder'' o
risco em pacotes complexos.

\section{Capítulo 2: O Vácuo das Pontocom e o 11 de Setembro
(2000-2001)}\label{capuxedtulo-2-o-vuxe1cuo-das-pontocom-e-o-11-de-setembro-2000-2001}

O medo de uma depressão econômica faz o Fed pisar no acelerador. *
\textbf{Liquidez:} A taxa de juros (\textbf{Fed Funds Rate}) cai de
6,5\% para 1\%. * \textbf{O Desvio:} O capital foge do Nasdaq (quebrado)
e busca abrigo no tijolo: o setor imobiliário vira o novo porto seguro.

\section{Capítulo 3: O Contrato com o Diabo
(2004)}\label{capuxedtulo-3-o-contrato-com-o-diabo-2004}

A política habitacional encontra a ganância de Wall Street. *
\textbf{GSE Policy Shift:} Fannie Mae e Freddie Mac recebem metas
agressivas para comprar hipotecas de baixa renda. * \textbf{A Banca do
João:} Bancos locais param de se preocupar com a qualidade do crédito
(Empréstimos NINJA) porque sabem que podem repassar o risco em 30 dias.

\section{Capítulo 4: A Alquimia Financeira e o Mercado
OTC}\label{capuxedtulo-4-a-alquimia-financeira-e-o-mercado-otc}

Como o lixo foi vendido como ouro. * \textbf{Securitização:} Hipotecas
podres são fatiadas e misturadas em CDOs. * \textbf{O Ponto Cego:} Como
não havia bolsa (Mercado de Balcão/OTC), os preços eram baseados em
modelos, não na realidade. * \textbf{Agências de Rating:} O carimbo
``AAA'' deu aos fundos de pensão a permissão legal para comprar risco
sistêmico.

\section{Capítulo 5: O Margin Call Global
(2007-2008)}\label{capuxedtulo-5-o-margin-call-global-2007-2008}

O castelo de cartas começa a balançar com a subida dos juros. *
\textbf{O Gatilho:} O Fed sobe juros para 5,25\% \(\rightarrow\) As
parcelas do ``Daniel'' dobram \(\rightarrow\) A inadimplência explode. *
\textbf{15/09/2008:} Lehman Brothers declara falência. O mercado
interbancário trava (ninguém confia em ninguém). * \textbf{16/09/2008:}
O resgate da \textbf{AIG} prova que o sistema estava segurado por uma
empresa que não tinha dinheiro para cobrir o sinistro.

\section{}\label{section}




\end{document}
